\documentclass[11pt,a4paper]{article}

% =============================
% PACKAGES
% =============================
\usepackage[margin=2.5cm]{geometry}
\usepackage{graphicx}      % For figures
\usepackage{amsmath, amssymb}  % Math symbols
\usepackage{booktabs}      % Nice tables
\usepackage{float}         % Control float position
\usepackage{caption}
\usepackage{subcaption}
\usepackage{hyperref}      % Clickable references
\usepackage{xcolor}        % Colors
\usepackage{enumitem}      % Better lists
\usepackage{url}           % URLs in text
\usepackage{listings}      % Code listings
\usepackage{appendix}      % Appendices
\usepackage{hyperref}
\usepackage{csvsimple}
\usepackage{xcolor}
\usepackage{makecell}

\definecolor{codegreen}{rgb}{0,0.6,0}
\definecolor{codegray}{rgb}{0.5,0.5,0.5}
\definecolor{codepurple}{rgb}{0.58,0,0.82}
\definecolor{backcolour}{rgb}{0.95,0.95,0.92}

\lstdefinestyle{mystyle}{
    backgroundcolor=\color{backcolour},   
    commentstyle=\color{codegreen},
    keywordstyle=\color{magenta},
    numberstyle=\tiny\color{codegray},
    stringstyle=\color{codepurple},
    basicstyle=\ttfamily\footnotesize,
    breakatwhitespace=false,         
    breaklines=true,                 
    captionpos=b,                    
    keepspaces=true,                 
    numbers=left,                    
    numbersep=5pt,                  
    showspaces=false,                
    showstringspaces=false,
    showtabs=false,                  
    tabsize=2
}

\lstset{style=mystyle}

% =============================
% TITLE INFORMATION
% =============================
\title{\textbf{CSP-Bender Problem}\\
\large Formulation, Solution in AMPL, and Comparison with Heuristic Methods}

\author{
  \textbf{Liliu Martinez} \and
  \textbf{Enrique Takiguchi} \\  
%   \small Master's in Statistics and Operations Research, IT2 Track \\
%   \small University of [Your University] \\
%   \small Date: \today
}

\date{}

% =============================
% DOCUMENT
% =============================
\begin{document}

\maketitle

\tableofcontents
\newpage

% =============================
% 1. INTRODUCTION
% =============================
\section{Introduction}
\label{sec:intro}

The Cell Supression Problem (CSP) is a fundamental problem in tabular data protection,
as it must be solved to ensure a certain level of confidenciality is applied to specific
sensible cells. Tabular data with unsolved CSP takes the risk that an attacker might be 
able to compute an interval for the values of sensible censured cells which is too small.

Since supressing only sensible cells would not suffice, the CSP seeks to find a balance
between the amount of information and the protection of the tables, by defining minimum
protection levels and the cost of hidding each cell.

As we will see, the complexity of the raw problem requires an algorithm to solve the problem
in an efficient way, for which we will use Benders decomposition.


\textbf{Objectives:}
\begin{itemize}[noitemsep]
  \item Contextualizing the problem
  \item Formulating the problem
  \item Implementing CSP in its basic form (MILP) in AMPL
  \item Implementing CSP with Benders decomposition in AMPL
  \item Testing the implementations
  
\end{itemize}

% =============================
% 2. PROBLEM DESCRIPTION AND FORMULATION
% =============================
\section{Problem Description and Formulation}
\label{sec:prob-descr}

\subsection{Basic definitions}
\begin{itemize}
\item \textbf{Microdata}: a file $V$ of $s$ individuals and $t$ variables, as a
matrix $s \times t$ where $V_{ij}$ is the value of variable $j$ for individual $i$.
\item \textbf{Tabular data}: tables generated crossing one ore more categorical variables 
from microdata.
These can be multidimensional tables holding frequency or magnitude information.\\
The table will have $n$ cells which will have to verify $m$ linear constraints.\\
The cells in the table are bounded by an upper and lower limit: $u<a_i<l$
\item \textbf{Protection levels from the p\% rule}:

An attacker might estimate the intervals of a cell
$[\overline{a}_p,  \underbar{$a$}_p]$, to assert these are not informative enough,
we must define some lower and upper protection level bounds to ensure this interval
fully contains the interval
\begin{equation}
  \label{eq:prot-lvls}
  [\overline{a}_p,  \underbar{$a$}_p]\in[a_p-lpl_p, a_p+upl_p]
\end{equation}

The best guess an auditor can make about $a_p$ is defined by $[\overline{a}_p,  \underbar{$a$}_p]\rightarrow$

\begin{equation}
\underline{a}_p = 
\begin{array}{ll} 
\min & x_p \\ 
\text{s.to} & Ax = b \\ 
& l_i \leq x_i \leq u_i \quad i \in \mathcal{P} \cup \mathcal{S} \\ 
& x_i = a_i \quad i \notin \mathcal{P} \cup \mathcal{S} 
\end{array} 
\qquad 
\overline{a}_p = 
\begin{array}{ll} 
\max & x_p \\ 
\text{s.to} & Ax = b \\ 
& l_i \leq x_i \leq u_i \quad i \in \mathcal{P} \cup \mathcal{S} \\ 
& x_i = a_i \quad i \notin \mathcal{P} \cup \mathcal{S} 
\end{array} 
\end{equation}


\end{itemize}



 
\subsection{CSP Formulation}

The Cell Suppression Problem seeks to hide sensitive (primary, $\mathcal{P}$) cells while minimizing
the suppression cost. Hiding only primary cells would not be enough since we could still infer
information about them from the rest of cells, thus, a number of secondary cells ($\mathcal{S}$) will
also be suppresed.

CSP computes the set $\mathcal{S}$ of secondary cells to be suppressed such that for all $p\in\mathcal{P}$, 
Equation \ref{eq:prot-lvls} is satisfied, where $\overline{a}_p,  \underbar{$a$}_p$ can be obtained with a minimization problem.

The CSP problem can be formulated as a Mixed-Integer Linear Programming (MILP) problem.
The MILP will be a minimization problem in terms of the following variables
\begin{itemize}
  \item $y_i\in\{0,1\}, i=1, ...n$: Indicates if a cell $i$ is suppressed (1) or not (0)
  \item $x^{l,p}\in \mathbb{R}^n$: Lower deviation of the primary cell $p$ from its original value
  \item $x^{u,p}\in \mathbb{R}^n$: Upper deviation of the primary cell $p$ from its original value
\end{itemize}

MILP formulation for CSP:
\begin{align*}
\min \quad & \sum_{i=1}^{n} w_i y_i \\
\text{s.to} \quad & \\
& \left. \begin{array}{rll}
(l_i - a_i)y_i \leq & \begin{array}{l} Ax^{l,p} \\ x_i^{l,p} \\ x_p^{l,p} \end{array} & \begin{array}{l} = 0 \\ \leq (u_i - a_i)y_i \\ \leq -lpl_p \end{array} \quad i = 1, \dots, n \\ \\
(l_i - a_i)y_i \leq & \begin{array}{l} Ax^{u,p} \\ x_i^{u,p} \\ x_p^{u,p} \end{array} & \begin{array}{l} = 0 \\ \leq (u_i - a_i)y_i \\ \geq upl_p \end{array} \quad i = 1, \dots, n
\end{array} \right\} \forall p \in \mathcal{P} \\
& y_i \in \{0, 1\} \quad x\ge 0\quad i = 1, \dots, n,
\end{align*}

The problem with the MILP formulation is that it becomes infeasibly complex as the table grows large.

The problem has:
\begin{itemize}
  \item $n$ binary variables and $2n\mathcal{P}$ continuous variables
  \item $2(m+2n)|\mathcal{P}|$ constraints
\end{itemize}

\subsection{Benders decomposition} \label{sec:benders-decomposition-formulation}

Benders decomposition partitions the problem into a master and subproblems, and
iterates efficiently between them to find the optimal solution.
\begin{itemize}
  \item Master problem: focuses on the binary decisions ($y_i$), which determine whether a cell is published ($y_i=0$) or suppressed ($y_i=1$)
  \item Subproblems: verify if the current suppression pattern chosen by the Master Problem actually provides the required protection levels ($lpl$ and $upl$) for each primary sensitive cell.
\end{itemize}

\textbf{Original Problem}

Given the original problem: 
\begin{align*} 
(P) = \quad \min \quad & c^T x + q^T y \\
\text{s. to} \quad & Tx + Wy = h \\
& x \in X \\
& y \geq 0
\end{align*}    


The two following problems are derivated:
\begin{itemize}
  \item \textbf{Master Formulation}
  \begin{align*}
  (BP_r) = \quad \min \quad & z \\
  \text{s. to} \quad & z \geq c^T x + u^{i\top} (h - Tx) &  i \in I \subseteq \{1 \dots p\} \\
  & v^{j\top} (h - Tx) \leq 0 & j \in J \subseteq \{1 \dots q\} \\
  & x \in X.
  \end{align*}
  \item \textbf{Subproblem Formulation}

    \begin{align*}
  (Q_D) = \quad \max \quad & u^T(h - Tx) \\
  \text{s. to} \quad & W^T u \leq q \\
  & u \in \mathbb{R}^m
  \end{align*}
\end{itemize}

To take advantage of Benders decomposition we first reformulate the original MILP
problem to an equivalent of Benders starting problem.

\begin{align*}
\min \quad & \sum_{i=1}^{n} c_i y_i \\
\text{s. to} \quad & Ax^{l,p} = 0 & \forall p \in \mathcal{P} \\
& Ax^{u,p} = 0 & \forall p \in \mathcal{P} \\
& (l_i - a_i)y_i \leq x_i^{l,p} & i=1 \dots n, \forall p \in \mathcal{P} \\
& x_i^{l,p} \leq (u_i - a_i)y_i & i=1 \dots n, \forall p \in \mathcal{P} \\
& (l_i - a_i)y_i \leq x_i^{u,p} & i=1 \dots n, \forall p \in \mathcal{P} \\
& x_i^{u,p} \leq (u_i - a_i)y_i & i=1 \dots n, \forall p \in \mathcal{P} \\
& x_p^{l,p} \leq -lpl_p & \forall p \in \mathcal{P} \\
& x_p^{u,p} \geq upl_p & \forall p \in \mathcal{P} \\
& y_i \in \{0, 1\}, \quad x, s \geq 0
\end{align*}

Given that $y_i$ belongs to a restricted set and both $x$ and $s$ are only restricted to be positive,
we take $y_i$ as the $x$ in Benders and $[x,s]$ as the $y$.

The Original Problem now becomes:
\begin{table}[h]
    \centering
    \begin{tabular}{| l | c c |}
        \toprule
        & \textbf{c\textsuperscript{T}x} & \textbf{q\textsuperscript{T}y} \\
        \hline
        \textbf{Objective function (min)} & $\sum_{i=1}^{n} c_i y_i$ & $+ 0 \sum_{i=1}^{n}x_{i}^{l,p}+x_{i}^{u, p}$ \\
        \bottomrule
    \end{tabular}
\end{table}

\begin{table}[h]
    \centering
    \begin{tabular}{| l | c c c |}
        \toprule
        \textbf{Constraint Type} & \textbf{Tx} & \textbf{+Wy} & \textbf{=h}\\
        \hline
        \textbf{Table Conservation} & 0 & \makecell[l]{$+ Ax^{l,p}$\\ $+ Ax^{u,p}$} & 0 \\
        \hline
        \textbf{Lower Capacity} & \makecell[l]{$(l - a)y$\\ $(l - a)y$} &  \makecell[l]{$-x^{l,p}$\\ $-x^{u,p}$} & $\le 0$ \\
        \hline
        \textbf{Upper Capacity} & \makecell[l]{$(a - u)y$\\ $(a - u)y$} & \makecell[l]{$+ x^{l,p}$ \\ $+ x^{u,p}$} & $\le 0$ \\
        \hline
        \textbf{Protection Levels} & 0 & \makecell[l]{$x^{l,p}$\\ $x^{u,p}$} & \makecell[l]{ $\le -lpl$\\ $ \ge upl$}\\
        \bottomrule
    \end{tabular}
\end{table}

\begin{table}[h]
\centering
\begin{tabular}{ |c | c c |}
    \hline
   & \textbf{$x \in X$} & \textbf{$y \geq 0$} \\
    \hline
  \textbf{Domain Constraints} & $y \in \{0, 1\}$ & $\quad x^{l}, x^{u} \geq 0$ \\
    \hline
\end{tabular}
\end{table}

\subsubsection{Subproblem} \label{sec:subproblem-formulation}
With this formulation we can now derive the Subproblem and Master Problem for Benders decomposition.
The subproblem will have to find the maximum deviation that an attacker can achieve for each primary cell $p$ while respecting the constraints of the table and it is whithin the protection levels.
Primal subproblem Q for each primary cell $p$:
\begin{itemize}
    \item Lower protection level subproblem:
        \begin{align*}
        (Q_) = \quad \min \quad & x_{p} \\
        \text{s. to} \quad & Ax = b [\lambda]\\
        & (l_i - a_i)\bar{y}_i \leq x_{i} [\mu^{l}]\\
        &  x_{i} \leq (u_i - a_i)\bar{y}_i [\mu^{u}]\\
        & -x_{p} \leq  lpl_{p} [\gamma]\\
        \end{align*}


    \item  Upper protection level subproblem:
        \begin{align*}
        (Q_) = \quad \min \quad & x_{p[j]} \\
        \text{s. to} \quad & Ax = b [\lambda]\\
        & (l_i - a_i)\bar{y}_i \leq x_{i} [\mu^{l}]\\
        &  x_{i} \leq (u_i - a_i)\bar{y}_i [\mu^{u}]\\
        & x_{p[j]} \leq  upl_{p[j]}  [\gamma]\\ 
        \end{align*}
\end{itemize}


Therefore, if we denote $protection\_lvl$ as either $lpl_{p}$ or $upl_{p}$, and $L$ as the corresponding correct sign, both subproblems can be summarized as:
\begin{align*}
(Q_) = \quad \min \quad & x_{p} \\
\text{s. to} \quad & Ax = b [\lambda]\\
& (l_i - a_i)\bar{y}_i \leq x_{i} [\mu^{l}]\\
&  x_{i} \leq (u_i - a_i)\bar{y}_i [\mu^{u}]\\
& L x_{p} \leq  protection\_lvl_{p}\\
\end{align*}

Given the primary, with $\mathcal{C}$ as the number of constraints in A matrix, we can derive the dual subproblem.
\begin{align*}
(Q_{D}) = \quad \max \quad & \sum_{c \in \mathcal{C}} \lambda_c b_c
+ \sum_{i=1}^{n}((l_i - a_i)\mu_{i}^{l} + (u_i-a_i)\mu_{i}^{u})\bar{y}_i  
+ \gamma protection\_lvl_{p[j]}\\
\text{s. to}  \quad & A^T\lambda + \mu^{l}  - \mu^{u} = e_p \gamma & i= 1..n \\
& \mu^{l} \geq 0, \mu^{u} \leq 0, \lambda free, \gamma \geq 0 
\end{align*}

\subsubsection{Benders Cuts}

With this changes, we no longer have lower and upper subproblems that gives us different cuts to add to the Master Problem.
The cuts to add to the Master Problem will depend on the optimal value of the combined dual subproblem objective such as:
\begin{itemize}
    \item Original bender cuts:
    \begin{itemize}
        \item Lower Benders cut:
            \begin{align*}
            -lpl \geq \quad & \sum_{i=1}^{n}((l_i - a_i)\mu_{ll,i} - (u_i-a_i)\mu_{ul,i})y_i  \\
            \end{align*}
        \item Upper Benders cut:
            \begin{align*}
            upl \leq \quad & \sum_{i=1}^{n}(-(l_i - a_i)\mu_{lu,i} + (u_i-a_i)\mu_{uu,i})y_i  \\
            \end{align*}
    \end{itemize}
    \item New Bender cuts:
        \begin{align*}
        \sum_{i=1}^{n}((l_i - a_i)\mu_{i}^{l} + (u_i-a_i)\mu_{i}^{u})y_i & \quad \forall k \in \mathcal{CUTS} \\
        + \gamma_{k} protection\_lvl_{k} \leq 0 \\
        \end{align*}
\end{itemize}

\subsubsection{Master Problem}
The Master Problem will now seek to minimize the suppression cost while respecting the Benders cuts added from previous iterations.

Given:
\begin{itemize}
  \item $\mathcal{CUTS}$: set of cuts added until iteration $r$ from insecure primary cells found with the lower dual subproblem
  \item for each cut $k \in \mathcal{CUTS}$, we have $\bar{\mu}_{l}^{p, k}, \bar{\mu}_{u}^{p, k}, \bar{\gamma}$ from the solution of the dual subproblem
\end{itemize}
The Master Problem at iteration $r$ is:
\begin{align*}
(BP_r) = \quad \min \quad & \sum_{i=1}^{n}c_i y_i \\
\text{s. to} \quad & y_p = 1 & \forall p \in \mathcal{P}\\
& \sum_{i=1}^{n}((l_i - a_i)\bar{\mu}_{i}^{l} + (u_i-a_i) \bar{\mu}_{i}^{u})y_i & \quad \forall k \in \mathcal{CUTS} \\
&  + \bar{\gamma}_{k} protection\_lvl_{k} \leq 0 \\
& y_i \in \{0, 1\}  
\end{align*}



% =============================
% 3. DATASETS
% =============================
\section{Datasets}
\label{sec:datasets}

\subsection{small}
Original table: 
In Table \ref{tab:small-original} we can see the original data table for \textit{small.ampl} data instance created by rows($\rightarrow$), with the sensitive cells marked with *.
Column 7 and Row 5 contain the totals for the rows and columns respectively.
\begin{table}[h]
    \centering
    \begin{tabular}{l | c c c c c c c}
             & Col 1 & Col 2 & Col 3 & Col 4 & Col 5 & Col 6 & Col 7 \\
        \toprule
        Row 1 & 300* & 8    & 5    & 5    & 5*   & 38   & \textbf{361} \\
        Row 2 & 8    & 68*  & 40   & 76   & 29   & 31   & \textbf{252} \\
        Row 3 & 11   & 33   & 20   & 60   & 35   & 44   & \textbf{203} \\
        Row 4 & 7    & 28   & 36*  & 41   & 22   & 63   & \textbf{197} \\
        Row 5 & \textbf{326} & \textbf{137} & \textbf{101} & \textbf{182} & \textbf{91} & \textbf{176} & \textbf{1013} \\
    \end{tabular}
    \caption{Original Small Instance Data Table with Sensitive Cells Marked with *}
    \label{tab:small-original}
\end{table}

\begin{table}[h]
    \centering
    \begin{tabular}{c | c | c}
        \textbf{Number Cell} & \textbf{Lower Protection Level} & \textbf{Upper Protection Level} \\
        \toprule
        1	& 15  & 20 \\
        5	& 2   & 4 \\
        9	& 10  & 6 \\
       24	& 9   & 3 \\
    \end{tabular}
    \caption{Table of Small Instance Protection Levels for Sensitive Cells}
    \label{tab:small-protection-levels}
\end{table}

Table Constraints:
\begin{enumerate}
  \item[] \hspace*{-\leftmargin}\textbf{Row Constraints:}
  \item  $x_{1} + x_{2} + x_{3} + x_{4} + x_{5} + x_{6} - x_{7} = 0$
  \item  $x_{8} + x_{9} + x_{10} + x_{11} + x_{12} + x_{13} - x_{14} = 0$
  \item  $x_{15} + x_{16} + x_{17} + x_{18} + x_{19} + x_{20} - x_{21} = 0$
  \item  $x_{22} + x_{23} + x_{24} + x_{25} + x_{26} + x_{27} - x_{28} = 0$
  \item[] \hspace*{-\leftmargin}\textbf{Column Constraints:}
  \item  $x_{1} + x_{8} + x_{15} + x_{22} - x_{29} = 0$
  \item  $x_{2} + x_{9} + x_{16} + x_{23} - x_{30} = 0$
  \item  $x_{3} + x_{10} + x_{17} + x_{24} - x_{31} = 0$
  \item  $x_{4} + x_{11} + x_{18} + x_{25} - x_{32} = 0$
  \item  $x_{5} + x_{12} + x_{19} + x_{26} - x_{33} = 0$
  \item  $x_{6} + x_{13} + x_{20} + x_{27} - x_{34} = 0$
\end{enumerate}


\subsection{example\_2D}

In Table \ref{tab:example2D-original} we can see the original data table for \textit{example\_2D.ampl} data instance created by rows($\rightarrow$), with the sensitive cells marked with *.
Column 1 and Row 1 represent the totals for their respective dimensions.
\begin{table}[h]
    \centering
    \begin{tabular}{l | c c c c c c}
         & Col 1 & Col 2 & Col 3 & Col 4 & Col 5 & Col 6 \\
        \toprule
        Row 1     & \textbf{3220} & \textbf{84} & \textbf{632} & \textbf{930} & \textbf{843} & \textbf{731} \\
        Row 2     & \textbf{1529} & 3  & 336 & 309  & 484 & 397 \\
        Row 3     & \textbf{484}  & 25 & 3   & 393* & 48  & 15  \\
        Row 4     & \textbf{392}  & 1  & 2   & 137* & 145 & 107  \\
        Row 5     & \textbf{815}  & 55 & 291* & 91  & 166 & 212* \\
    \end{tabular}
    \caption{Original Example\_2D Instance Data Table with Sensitive Cells Marked with *}
    \label{tab:example2D-original}
\end{table}

\begin{table}[h]
    \centering
    \begin{tabular}{c | c | c}
        \textbf{Number Cell} & \textbf{Lower Protection Level} & \textbf{Upper Protection Level} \\
        \toprule
        16	&	40	&	30 \\
	      22	&	14	&	14 \\
	      27	&	15	&	30 \\
	      30	&	21	&	21
    \end{tabular}
    \caption{Table of Example\_2D Instance Protection Levels for Sensitive Cells}
    \label{tab:example2D-protection-levels}
\end{table}


\subsection{targus}
In Table \ref{tab:targus-original} we can see the original data table for \textit{targus.ampl} data instance created by rows($\rightarrow$), with the sensitive cells marked with *. 
Due to the large values in the table, we have used commas to separate thousands.
Due to the row-wise size of the table, we have transposed the original table for better visualization.
Row 1 contains the totals for the columns.
\begin{table}[H]
    \centering
    \begin{tabular}{l | c c c c c c c c c}
         & Row 1 & Row 2 & Row 3 & Row 4 & Row 5 & Row 6 & Row 7 & Row 8 & Row 9\\
        \toprule
        Col 1 & 16847261.84 & 20* & 25 & 2711808 & 2320534 & 2505042.58 & 2799074.26 & 6510758 & 0 \\
        Col 2 & 4373279     & 5* & 5* & 719049 & 659680 & 688962 & 756529 & 1549049 & 0 \\
        Col 3 & 1986129     & 5* & 5* & 398062 & 348039 & 354711 & 418778 & 466529 & 0 \\
        Col 4 & 1808861     & 0 & 0 & 223990 & 221332 & 241913 & 258233 & 863393 & 0 \\
        Col 5 & 578289      & 0 & 0 & 96997 & 90309 & 92338 & 79518 & 219127 & 0 \\
        Col 6 & 3703896     & 15* & 5* & 642238 & 515003 & 534147 & 620392 & 1392096 & 0 \\
        Col 7 & 124336      & 5* & 0 & 36311 & 32132 & 25770 & 18150 & 11968* & 0 \\
        Col 8 & 526279      & 0 & 0 & 93589 & 94957 & 110930 & 81799 & 145004 & 0 \\
        Col 9 & 2234995     & 10* & 5* & 345803 & 251358 & 251188 & 303377 & 1083254 & 0 \\
        Col 10 & 818286     & 0 & 0 & 166535 & 136556 & 146259 & 217066 & 151870 & 0 \\
        Col 11 & 4576115.84 & 0 & 0 & 648972 & 543570 & 663896.58 & 775132.26 & 1944545 & 0 \\
        Col 12 & 485326     & 0 & 0 & 63767 & 75442 & 87305 & 59953 & 198859 & 0 \\
        Col 13 & 3664559.84 & 0 & 0 & 537911 & 430851 & 515019.58 & 643762.26 & 1537016 & 0 \\
        Col 14 & 426230     & 0 & 0 & 47294 & 37277 & 61572 & 71417 & 208670 & 0 \\
        Col 15 & 4193971    & 0 & 15* & 701549 & 602281 & 618037 & 647021 & 1625068 & 0 \\
        Col 16 & 2752743    & 0 & 15* & 488613 & 392395 & 363490 & 402925 & 1105305 & 0 \\
        Col 17 & 1441228    & 0 & 0 & 212936 & 209886 & 254547 & 244096 & 519763 & 0 \\
        Col 18 & 0          & 0 & 0 & 0 & 0 & 0 & 0 & 0 & 0 \\
      \end{tabular}
    \caption{Original Targus Instance Data Table with Sensitive Cells Marked with *}
    \label{tab:targus-original}
\end{table}


\begin{table}[H]
    \centering
    \begin{tabular}{c | c | c}
        \textbf{Number Cell} & \textbf{Lower Protection Level} & \textbf{Upper Protection Level} \\
        \toprule
       19	& 6   & 6 \\
       20	& 1.67 & 1.67 \\
       21	& 1.67 & 1.67 \\
       24	& 5   & 5 \\
       25	& 1.67 & 1.67 \\
       27	& 3.33 & 3.33 \\
       38	& 1.67 & 1.67 \\
       39	& 1.67 & 1.67 \\
       42	& 1.67 & 1.67 \\
       45	& 1.67 & 1.67 \\
       51	& 5   & 5 \\
       52	& 5   & 5 \\
       133	& 3480 & 3480
    \end{tabular}
    \caption{Table of Targus Instance Protection Levels for Sensitive Cells}
    \label{tab:targus-protection-levels}
\end{table}



\section{Implementation Details}
Given the the mathematical formulation of the Benders decomposition from Section \ref{sec:benders-decomposition-formulation}, we have implemented the master problem and subproblems:
\subsection{Subproblem Lower}


\subsection{Subproblem Upper}


\section{Computational Experiments and Results}


\subsection{small}

\begin{table}[H]
    \centering
    \begin{tabular}{c | c | c}
        \textbf{Cell ID} & \textbf{Cost} & \textbf{Is Primary} \\
        \toprule
        1          & 1          & 1   \\      
        3          & 1          & 0   \\    
        5          & 1          & 1   \\     
        8          & 1          & 0   \\    
        9          & 1          & 1   \\    
        10         & 1          & 0   \\    
        12         & 1          & 0   \\    
        15         & 1          & 0   \\    
        22         & 1          & 0   \\    
        24         & 1          & 1   \\
    \end{tabular}
    \caption{Table of Cells to be Suppressed in Small Instance}
    \label{tab:small-suppressed-cells}
\end{table}


\begin{table}[H]
    \centering
    \begin{tabular}{l | c c c c c c c}
             & Col 1 & Col 2 & Col 3 & Col 4 & Col 5 & Col 6 & Col 7 \\
        \toprule
        Row 1 & x    & 8    & x    & 5    & x   & 38   & \textbf{361} \\
        Row 2 & x    & x    & x    & 76   & x   & 31   & \textbf{252} \\
        Row 3 & x    & 33   & 20   & 60   & 35   & 44   & \textbf{203} \\
        Row 4 & x    & 28   & x    & 41   & 22   & 63   & \textbf{197} \\
        Row 5 & \textbf{326} & \textbf{137} & \textbf{101} & \textbf{182} & \textbf{91} & \textbf{176} & \textbf{1013} \\
    \end{tabular}
    \caption{Small Data Table with Suppressing Sensitive Cells}
    \label{tab:small-suppressed-table}
\end{table}





\subsection{Example\_2D}

\begin{table}[H]
    \centering
    \begin{tabular}{c | c | c}
        \textbf{Cell ID} & \textbf{Cost} & \textbf{Is Primary} \\
        \toprule
        3          & 0          & 0 \\        
        6          & 0          & 0 \\      
        13         & 0          & 0 \\      
        16         & 0          & 1 \\      
        19         & 0          & 0 \\      
        22         & 0          & 1 \\      
        27         & 0          & 1 \\      
        30         & 0          & 1 \\
    \end{tabular}
    \caption{Table of Cells to be Suppressed in example\_2D Instance}
    \label{tab:example2D-suppressed-cells}
\end{table}


\begin{table}[H]
    \centering
    \begin{tabular}{l | c c c c c c}
         & Col 1 & Col 2 & Col 3 & Col 4 & Col 5 & Col 6 \\
        \toprule
        Row 1     & \textbf{3220} & \textbf{84} & \textbf{x} & \textbf{930} & \textbf{843} & \textbf{x} \\
        Row 2     & \textbf{1529} & 3  & 336 & 309  & 484 & 397 \\
        Row 3     & \textbf{x}    & 25 & 3   & x & 48  & 15  \\
        Row 4     & \textbf{x}    & 1  & 2   & x & 145 & 107  \\
        Row 5     & \textbf{815}  & 55 & x   & 91  & 166 & x \\
    \end{tabular}
    \caption{example\_2D Data Table with Suppressing Sensitive Cells}
    \label{tab:example2D-suppressed-table}
\end{table}




\subsection{targus}

\begin{table}[H]
    \centering
    \begin{tabular}{c | c | c}
        \textbf{Cell ID} & \textbf{Cost} & \textbf{Is Primary} \\
        \toprule
        19         & 1          & 1          \\
        20         & 1          & 1         \\
        21         & 1          & 1         \\
        24         & 1          & 1         \\
        25         & 1          & 1         \\
        27         & 1          & 1         \\
        37         & 1          & 0         \\
        38         & 1          & 1         \\
        39         & 1          & 1         \\
        42         & 1          & 1         \\
        45         & 1          & 1         \\
        51         & 1          & 1         \\
        52         & 1          & 1         \\
        78         & 611        & 0         \\
        87         & 715        & 0         \\
        99         & 298        & 0         \\
        106        & 432        & 0         \\
        133        & 14         & 1         \\
        134        & 172        & 0  \\
    \end{tabular}
    \caption{Table of Cells to be Suppressed in Targus Instance}
    \label{tab:targus-suppressed-cells}
\end{table}

\begin{table}[H]
    \centering
    \begin{tabular}{l | c c c c c c c c c}
         & Row 1 & Row 2 & Row 3 & Row 4 & Row 5 & Row 6 & Row 7 & Row 8 & Row 9\\
        \toprule
        Col 1 & 16847261.84 & x & x & 2711808 & 2320534 & 2505042.58 & 2799074.26 & 6510758 & 0 \\
        Col 2 & 4373279     & x & x & 719049 & 659680 & 688962 & 756529 & 1549049 & 0 \\
        Col 3 & 1986129     & x & x & 398062 & 348039 & 354711 & 418778 & 466529 & 0 \\
        Col 4 & 1808861     & 0 & 0 & 223990 & 221332 & 241913 & 258233 & 863393 & 0 \\
        Col 5 & 578289      & 0 & 0 & 96997 & 90309 & 92338 & 79518 & 219127 & 0 \\
        Col 6 & 3703896     & x & x & 642238 & x     & 534147 & 620392 & 1392096 & 0 \\
        Col 7 & 124336      & x & 0 & 36311 & 32132 & 25770 & 18150 & x & 0 \\
        Col 8 & 526279      & 0 & 0 & 93589 & 94957 & 110930 & 81799 & x & 0 \\
        Col 9 & 2234995     & x & x & 345803 & 251358 & x & 303377 & 1083254 & 0 \\
        Col 10 & 818286     & 0 & 0 & 166535 & 136556 & 146259 & 217066 & 151870 & 0 \\
        Col 11 & 4576115.84 & 0 & 0 & 648972 & 543570 & 663896.58 & 775132.26 & 1944545 & 0 \\
        Col 12 & 485326     & 0 & 0 & 63767 & 75442 & 87305 & 59953 & 198859 & 0 \\
        Col 13 & 3664559.84 & 0 & 0 & 537911 & 430851 & 515019.58 & 643762.26 & 1537016 & 0 \\
        Col 14 & 426230     & 0 & 0 & 47294 & 37277 & 61572 & 71417 & 208670 & 0 \\
        Col 15 & 4193971    & 0 & x & 701549 & x    & 618037 & 647021 & 1625068 & 0 \\
        Col 16 & 2752743    & 0 & x & 488613 & 392395 & x & 402925 & 1105305 & 0 \\
        Col 17 & 1441228    & 0 & 0 & 212936 & 209886 & 254547 & 244096 & 519763 & 0 \\
        Col 18 & 0          & 0 & 0 & 0 & 0 & 0 & 0 & 0 & 0 \\
      \end{tabular}
    \caption{Targus Data Table with Suppressing Sensitive Cells}
    \label{tab:targus-suppressed-table}
\end{table}

% =============================
% 9. REFERENCES
% =============================
\section{References}

% load bibliography file
% \bibliographystyle{plain}
% \bibliography{citations}



% \begin{itemize}
%   \item Jain, A. K., and Dubes, R. C. (1988). \textit{Algorithms for Clustering Data}. Prentice Hall.
%   \item Kaufman, L., and Rousseeuw, P. J. (1990). \textit{Finding Groups in Data: An Introduction to Cluster Analysis}.
%   \item UCI Machine Learning Repository: \url{https://archive.ics.uci.edu/}
%   \item Tocto Inga, P. M. (2012). \textit{Comparación de los algoritmos Prim y Kruskal [Tesis de licenciatura, Universidad Nacional de Ingeniería, Facultad de Ciencias]}. Universidad Nacional de Ingeniería.
% \end{itemize}

\newpage


% =============================
% APPENDIX
% =============================



% Listing of AMPL files
\begin{appendix}
\section{AMPL Files}
\label{apdx:ampl-files}
\begin{lstlisting}[language=Python, caption=\texttt{kmedoid.mod}]
param n; # observations
param k; # clusters
set M:={1..n};

param d{M,M};
var x{M,M} binary;


minimize f: sum{i in M}
    sum{j in M} d[i,j]*x[i,j];

subject to one_cluster{i in M}:
    sum{j in M} x[i,j] = 1;

subject to k_clusters:
    sum{j in M} x[j,j] = k;

subject to cluster_belongs{i in M, j in M}:
    x[j,j] >= x[i,j];
\end{lstlisting}

\begin{lstlisting}[language=Python, caption=\texttt{kmedoid.run}]
# Reset enviroment
reset;

# Load model and data
model kmedoids.mod;
data kmedoids.dat;

option solver cplex;
solve;

display _solve_elapsed_time;

printf {i in M, j in M} "%s,%s,%i\n", i, j, x[i,j] > "x_matrix.csv";
\end{lstlisting}

\begin{lstlisting}[language=Python, caption=Python function to generate \texttt{kmedoid.dat} ]
def generate_dat(N, D, k, filename='AMPL/kmedoid.dat'):
    with open(filename, 'w') as f:
        f.write("#Define sets\n")
        f.write(f"param n := {N};\n")
        f.write(f"param k := {k};\n")

        f.write("#Define parameters\n")
        f.write(f"param d: {'\t'.join([str(j+1) for j in range(N)])} :=\n")
        for i in range(N):
            f.write(f" {i+1} " + ' '.join(f"{D[i][j]:.6f}" for j in range(N)) + "\n")
        f.write(";\n")

    \end{lstlisting}
  \end{appendix}

\end{document}
